\documentclass{jarticle}

\title{簡易木製スピニングリール制作について}

\begin{document}


\maketitle

\section{初めに}
この文章を通読する必要はない.
写真だけ,　設計図だけ見て作り始めるのも可.
試行錯誤もまた工作の楽しみの一つである.

\subsection{先に言っておきたいこと}
\begin{enumerate}
\item 制作過程を楽しむ(楽しめる)ことがわりあい大切.
\item 市販品の性能を求めては落胆する. ロマンを重視しよう.
\end{enumerate}

※余談になるが, 
リールをできるだけ安く手に入れたいという方向に効率を求める場合は
釣り人の知人を作ることを勧める.
趣味人というものは, まだまだ使える道具があっても, つい新しいものに手を出してしまいがちで,
性能も状態も十分な道具を余らせがちなものである.

\subsection{「簡易」木製スピニングリールとは}
設計するにあたって「スピニングリール」を
	回転するアームと, 竿と平行に固定された糸巻を持ち,
	次の機構を持っているものととらえている.
	\begin{enumerate}
	\item
		糸を回転するアームに引っ掛けて, 糸巻に巻きつけることによって巻き取る.
	\item 
		アームから糸を外すことにより, 糸の自由放出(とそれを利用した仕掛けの投擲)が可能.
	\end{enumerate}
	制作難度の調整結果としてペールやオシレーション機構,ドラグ機構等は省略している.
	\begin{enumerate}
	\item
		ペールの省略:糸の引っ掛けが手動なのでもたつく→がんばれ
	\item 
		オシレーション機構の省略:糸絡みの可能性の増大→そんなにヘビーに使うもんではない
	\item 
		ドラグ機構の省略:ロマンが減る→ロマン成分は自作に求める
	\end{enumerate}



\section{制作}
\subsection{材料}
\begin{enumerate}
\item 100mm x 100mm x 6mm MDF板, 4枚（DAISOで6枚組で売ってる）
\item 丸箸 (コンビニ等で弁当とかに付けてくれるやつ) (約6mmm径の丸棒)
\item 竹串 (約3mmm径の丸棒)
\item 木工ボンド
\item 瞬間接着剤
\end{enumerate}
2～5も家のどっかになければ100均でも入手可.
2，5はそれっぽい棒で用意したドリル刃の径と適合してれば何でも代用可.

実際に遣うためには以下も必要
\begin{enumerate}
\item 巻くための釣り糸(ナイロン糸で十分)
\item 竿取付用金具(適当な大きさのステー等で代用可, 何なら竿も自作してねじ止めしても良し)
\end{enumerate}

\subsection{道具}
\begin{enumerate}
\item のこぎり
\item 平彫刻刀あるいは小さなノミ（C2部材の溝切り落としに使うだけ）
\item ラチェットドライバーとドリル刃
	(6mmと2.8mm, 丸箸と竹串の直径に大体あってるもの.
	多少いい加減でも良いが太い方は緩く細い方はきつめの方が都合良い.)
\item コンパスと定規, 筆記具
\item (木工)やすり
\end{enumerate}

あると便利な道具
\begin{enumerate}
\item マスキングテープ（仮固定用）
\item クランプ（固定, 圧着用）
\item 丸やすり（穴径の調整）
\item 垂直穴あけ用のジグ
\item ホビーカンナ(円に削るときとか早くなる)
\item 止型定規（線引きに何かと便利）
\end{enumerate}

\subsection{作り方}
最重要部分はAパーツBパーツの2つの歯車部分で, 特に歯の部分の(穴位置)精度が重要なのでそこだけは気を付けること.
逆を言えばそれ以外の部分は多少適当でも動く.

\subsubsection{部材の切り出し}



円はのこぎりで大雑把に丸くした後, やすりなどで丸くする.
実を言えば機械機構としての円の作成を要求する部位はない(丸棒の丸とドリル穴の丸を使っている)のでそこまで頑張って丸くしなくても良い.
A2B1D1が丸いのは他の部材との干渉を避けるためであり, D2D3が丸いのは糸の出をスムーズにする為であるので
これらをクリアする程度の円でよい.

一方で, C123の長辺の直線, C2切り込みの長辺に対する直交性は機械機構として必要な要素.
長辺は切り出し前の直線を生かせると楽. 切り込みは丁寧に切ること.



\subsubsection{組み立て}
\begin{enumerate}
\item
	設計図に従ってABCD各パーツを組む.
	\begin{itemize}
	\item 基本木工ボンドで圧着
	\item AパーツBパーツの歯の根元に最後に瞬間接着剤をたらして強度を上げる
	\item A3パーツの糸の「入り口」はヤスリで削って少し広げておいた方が使い易い.
	\end{itemize}
\item
	AパーツにDパーツを挿し, Aパーツが回転できるようにしながらDパーツをCパーツに固定する.
	Dパーツは回転してはいけない.(引き抜く方向に力がかかるのでそれに抗するように固定)
\item
	適当な長さに切った丸箸をBパーツとCパーツに通し,
	回転可能な様に箸の両端を止める.
\end{enumerate}
各パーツどうしの固定は接着しない木釘(というか竹串釘)
を用いるなど取り外し可能な様にできればその方が使い勝手が良い.


\end{document}

